\section{A split \(\mathbb Q(\sqrt{-2})\) sixfold source from an \(F_4\) extension}
\label{sec:gf4-new}

Put \(K=\mathbb Q(t)\), \(t^2=-2\), and \(\mathcal O_K=\mathbb Z[t]\).
We construct a finite subgroup of \(U(4)\) over \(K\), a twelve-involution
cover, and an exact rank-nine period certificate.  Two local involutions
have a two-dimensional negative eigenspace; this is handled explicitly
in the parabolic calculation.  The resulting split sixfold source gives
an explicit effective detector bound.

\paragraph{Provenance.}
The standard \(F_4\) root system and four simple roots below are the
published ones in \cite[\S3, p.~83]{GF4:Cahn}.  The extra matrix, five
chosen reflection roots, and exact boundary certificate are constructed
and checked here.  The unrefereed rStation paper75 suggested an \(F_4\)
extension over \(\mathbb Q(\sqrt{-2})\); none of its theorem statements,
matrices, or numerical witnesses is used.  The proof establishes the
finite group directly rather than relying on an unverified classification
of possible sources.

\subsection*{D-GF4-01. The group and its integral sixfold factor}

Use the standard positive definite form \(\bar x^{\,t}y\) on \(K^4\).
For a nonzero vector \(v\), write
\[
                  s_v=1-2v\bar v^{\,t}/(\bar v^{\,t}v).
\]
Let \(R\) be the usual \(48\)-root system
\[
 R=\{\pm e_i\}\ \cup\ \{\pm e_i\pm e_j:i<j\}
       \ \cup\ \{(\pm e_1\pm e_2\pm e_3\pm e_4)/2\},
\]
and \(W=\langle s_v:v\in R\rangle=W(F_4)\).  It is generated by the
four simple reflections in
\[
 e_2-e_3,\quad e_3-e_4,\quad e_4,\quad(e_1-e_2-e_3-e_4)/2.
\]

Here is a direct order proof.  The signed permutation group \(B_4\)
lies in \(W\), and a reflection in a half-integral root sends a coordinate
root to a half-integral short root.  Thus \(W\) is transitive on the
\(24\) short roots.  Its stabilizer of \(e_1\) permutes the six short
roots perpendicular to \(e_1\), namely \(\pm e_2,\pm e_3,\pm e_4\).
Restriction embeds that stabilizer in the signed permutation group
\(B_3\), of order \(48\); the entire \(B_3\) already belongs to it.
Consequently \(|W|=24\cdot48=1152\).  The fact that the displayed
reflections preserve \(R\) follows directly from their formulas, and is
also checked by the finite reconstruction.

Define
\[
 \Phi=\begin{pmatrix}
 0&-1&-1&0\\1&0&0&1\\1&0&0&-1\\0&-1&1&0
 \end{pmatrix},\qquad z=\frac t2\Phi.
\]
We have \(\Phi^t=-\Phi\), \(\Phi^2=-2I\), and therefore
\(z^*=z\), \(z^2=1\).  Its trace is zero, so the negative and positive
eigenspaces both have dimension two.  The matrix \(\Phi\) sends every
short root to a long root and every long root to twice a short root,
up to sign.  For coordinate roots this follows from its columns.  For
half-integral roots its coordinates belong to \(\{0,1,-1\}\), and its
squared norm is two, so exactly two coordinates are nonzero.  The assertion
for long roots then follows from \(\Phi^2=-2I\).  Hence \(z\) permutes
the root lines and normalizes \(W\).  As \(z\) has nonreal entries,
\[
             G_F=\langle W,z\rangle=W\sqcup Wz,\qquad |G_F|=2304.
\]
This is a group over \(K\) preserving the standard positive form.

The subgroup \(B_4\) already makes the natural representation absolutely
irreducible: the independent coordinate sign changes force an invariant
complex subspace to be a sum of coordinate lines, and the permutations
act transitively on those lines.  Put \(s_0=s_{e_1}\), let \(P_{12}\)
swap the first two coordinates, and set \(g_0=P_{12}z\).  Directly,
\[
       (1-s_0)g_0e_1=te_1,
       \qquad \operatorname{tr}(s_0g_0)=-t.
\]
Thus its character field is exactly \(K\).  The displayed realization
has Schur index one, and its rational simple block is \(M_4(K)\).

Choose five reflection roots
\[
\begin{array}{c|rrrr}
 i&\multicolumn{4}{c}{v_i^t}\\\hline
1&0&-1&-1&0\\
2&0&-1&0&1\\
3&0&0&0&-1\\
4&-1/2&1/2&-1/2&-1/2\\
5&0&0&-1&1
\end{array}
\]
and set \(s_i=s_{v_i}\).  In terms of these, the four standard simple
reflections listed above are, respectively,
\[
                 s_2s_5s_2,\quad s_5,\quad s_3,\quad
                 s_1s_2s_4s_2s_1.
\]
These short word identities are checked by direct multiplication.  Thus
the \(s_i\) generate \(W\), and the doubled tuple
\[
       (s_1,s_1,s_2,s_2,s_3,s_3,s_4,s_4,s_5,s_5,z,z)
\]
has product one and generates \(G_F\).  All twelve inertia elements
are involutions.  Ten have negative-eigenspace dimension one, and two
have negative-eigenspace dimension two.

The connected cover \(C\to\mathbb P^1\) has
\[
 2g(C)-2=2304\bigl(-2+12(1-1/2)\bigr)=4\cdot2304,
 \qquad g(C)=4609.
\]
Chevalley--Weil \cite[Theorem 2.10]{G34:FGP} gives holomorphic multiplicity
\[
                     -4+\frac{10+2\cdot2}{2}=3
\]
for each conjugate natural character.  Its multiplicity factor is a
Weil sixfold of signature \((3,3)\).  Notice that the source genus uses
the order of each inertia element, whereas Chevalley--Weil uses its
negative-eigenspace multiplicity; the two rank-two involutions are not
being counted as reflections.

Let \(G=2304\), let \(\chi\) be the natural character, and define
\[
 \varepsilon=\frac4G\sum_{a\in G_F}
       \operatorname{Tr}_{K/\mathbb Q}\bigl(\chi(a^{-1})\bigr)a,
 \qquad e=\varepsilon(1-s_0)/2.
\]
Then \(e\) is a rational primitive Rosati-symmetric idempotent and
\(Ge\in\mathbb Z[G_F]\), because \(4/2=2\) and the character traces
are integers.  Let \(A\subset J(C)\) be its image.  The integral
endomorphism \((1-s_0)P_{12}z\) preserves \(A\) and restricts to \(t\)
on its natural root line, and to \(-t\) on the conjugate line.  It therefore
supplies the integral \(\mathbb Z[t]\)-action on \(A\).  As in the
dual-lattice argument of D-G24-01, \(G\Lambda_A^\#\subseteq\Lambda_A\).
For the restricted principal polarization,
\begin{equation}\label{eq:gf4-source-delta}
        \Delta_A\mid G^6,\qquad \int_A\theta_A^6\le6!G^6.
\end{equation}

\subsection*{D-GF4-02. Parabolic coordinates with rank-two inertia}

For the first five negative eigenspaces use the column vectors
\[
\begin{gathered}
 q_1=(0,1,1,0)^t,\quad q_2=(0,1,0,-1)^t,\quad
 q_3=(0,0,0,2)^t,\\
 q_4=(1,-1,1,1)^t/2,\qquad q_5=(0,0,1,-1)^t.
\end{gathered}
For \(\operatorname{im}(1-z)\) use
\[
                   q_6=(1,-t/2,-t/2,0)^t,
                   \qquad q_7=(t/2,1,0,t/2)^t.
\]
These are obtained by taking independent columns of \(1-t_i\), and
their asserted eigenspace dimensions are verified exactly.

Write the cocycle values at the first ten branch loops as
\(u_i=\lambda_iq_{\lceil i/2\rceil}\).  At the last two put
\[
          u_{11}=\lambda_{11}q_6+\lambda_{12}q_7,
          \qquad u_{12}=\lambda_{13}q_6+\lambda_{14}q_7.
\]
Thus the parabolic coefficient space has dimension fourteen, not twelve.
With \(Q_i=t_1\cdots t_{i-1}\), impose \(\sum_iQ_iu_i=0\), and quotient
by \(u_i=(1-t_i)a\).  The relation and coboundary matrices each have
rank four, their product is zero, and the quotient has dimension
\(14-4-4=6\).

The first three quotient representatives \(b_1,b_2,b_3\) have equal
coefficients one in the first, second and third branch pair respectively.
The remaining representatives are supported in coefficient positions
\((1,3,5,7,9,11,12)\), with rows
\[
\begin{aligned}
 b_4:\;&(-1,1,1,0,1,0,0),\\
 b_5:\;&(1+t/2,-2,-1/2,-2,0,1,0),\\
 b_6:\;&(t/2,-1-t,-1/2-t/2,-t,0,0,1).
\end{aligned}
\]
They satisfy the relation, and the six columns together with the four
coboundary columns have rank ten.  Hence they form a quotient basis.
Put \(I=\langle b_1,b_2,b_3\rangle_K\).

Define the auxiliary form
\[
 S(u,v)=\sum_i\left\langle Q_i^{-1}\sum_{j<i}Q_ju_j+u_i/2,v_i\right\rangle.
\]
Here \(u_i/2\) applies to the entire negative eigenspace, including the
two-dimensional one for \(z\); no rank-one formula is substituted there.
The exact computation gives a nondegenerate anti-Hermitian form with zero
upper-left \(3\)-by-\(3\) block.  Its full matrix is in the matrix appendix.

Use the same nine compatible intervals
\[
 [1,2],\ [3,4],\ [5,6],\ [7,8],\ [9,10],\ [11,12],
 \ [1,4],\ [5,8],\ [1,8].
\]
They form a maximal boundary stratum of \(\overline M_{0,12}\), and
every interval monodromy product equals one.  Applying the cocycle identity
to the returning boundary twist gives
\[
 u_i'=u_i+(1-t_i)\sum_{j=a}^b(t_a\cdots t_{j-1})u_j
       \quad(a\le i\le b),
\]
with unchanged outside coordinates.  Coordinates in each negative
eigenspace are solved in its displayed basis, then the quotient reduction
is performed using the verified ten-column basis.  The resulting matrices
satisfy
\[
 T_j=1+N_j,\qquad N_j=\begin{pmatrix}0&X_j\\0&0\end{pmatrix},
 \qquad N_iN_j=0,\qquad T_j^*ST_j=S.
\]
Their ranks, in interval order, are
\[
                         1,1,1,1,1,2,2,2,3.
\]

Minus the lower-right block of \(SN_j\) is Hermitian positive semidefinite.
Here are all its nonempty principal minors, ordered
\(1,2,3,12,13,23,123\):
\[
\begin{array}{c|rrrrrrr}
j&1&2&3&12&13&23&123\\\hline
1&2&3&1&0&0&0&0\\
2&2&8&6&0&0&0&0\\
3&4&1&3&0&0&0&0\\
4&0&4&2&0&0&0&0\\
5&2&0&0&0&0&0&0\\
6&0&2&2&0&0&2&0\\
7&2&7&5&9/2&9/2&3/2&0\\
8&4&7&7&12&6&9/2&0\\
9&2&2&2&7/2&7/2&2&2
\end{array}
\]
The leading principal minors of their sum are
\[
                             18,\quad379,\quad1286.
\]
Every positive weighted sum of the \(N_j\) thus has rank three, image
\(I\), and kernel \(I\).  The exact determinant obtained by flattening
the nine \(X_j\) row by row is
\begin{equation}\label{eq:gf4-rank-minor}
           \frac{256+640t}{27}\ne0,
           \qquad N_{K/\mathbb Q}\left(\frac{256+640t}{27}\right)
                  =\frac{2^{15}}{3^3}.
\end{equation}
This is a determinant over \(K\), and hence a nonzero determinant under
either complex embedding.

\subsection*{D-GF4-03. Actual period rank and splitness}

The geometric braid/cocycle comparison
\cite[Lemma 2.2 and Theorem 2.5]{G34:DWmon} applies to arbitrary finite
inertia, so it applies to the two rank-two involutions as well.  Finite
marking and a coordinatewise Kummer cover of the nine-node plumbing
polydisc give geometric logarithms
\[
                  M_j=\epsilon c_jN_j,
        \qquad c_j\in\mathbb Z_{>0},\quad \epsilon\in\{1,-1\},
\]
with one common sign from the orientation/transport convention.  The
finite central ambiguity is killed by the finite cover.  All products
\(M_iM_j\) vanish.  The positive-sum calculation makes
\(M=\sum_jM_j\) have rank three over \(K\) and image and kernel \(I\).

Write the period map as
\[
       F(w)=\exp\!\left(\sum_jw_jM_j\right)\Psi(z),
       \qquad z_j=e^{2\pi i w_j},\qquad \Psi(0)=F_0,
\]
using \cite[Theorem 4.12]{G34:Schmid}.  On the underlying rational space
the weight filtration of the diagonal limit is
\(W_0=W_1=I_{\mathbb Q}\), \(W_2=V_{\mathbb Q}\).
By \cite[Theorem 6.16]{G34:Schmid} it is a limiting mixed Hodge structure.
The conditions \(F_0^0=V_{\mathbb C}\), \(F_0^2=0\) imply that
\(\operatorname{Gr}_0^W\) has type \((0,0)\).  Therefore
\(F_0^1\cap I_{\mathbb C}=0\), and in either \(K\)-embedding the
three-dimensional \(F_0^1\) is a complement to \(I\).

At that point the Grassmannian tangent space is
\(\operatorname{Hom}(F_0^1,I)\).  The tangent vectors \(M_jF_0\)
are represented, up to fixed invertible changes of basis, by the nine
nonzero scalar multiples \(\epsilon c_jX_j\).  Equation
\eqref{eq:gf4-rank-minor} makes them independent.  Translating a period
derivative by \(\exp(-\sum_jw_jM_j)\) gives
\[
                  M_j\Psi(z)+2\pi i z_j\partial_j\Psi(z).
\]
Its limit is \(M_jF_0\), so the derivative determinant remains nonzero
at sufficiently nearby smooth points.  The actual period rank is nine,
and the signature-\((3,3)\) period image contains a nonempty analytic open.
This conclusion uses the interior derivative limit, not just the nine
boundary parameters or the dimension of their span over \(\mathbb Q\).

The genuine geometric logarithm \(M\) is skew-adjoint for the polarization
and has square zero.  Thus \(E(Mu,Mv)=0\).  Its \(K\)-rank-three image
is a maximal isotropic subspace, proving that the source Hermitian form
is split.  The auxiliary form was only used to certify the rank of every
positive weighted sum, not to identify the geometric polarization.

\subsection*{D-GF4-04. Suppliers and the explicit bound}

Every inertia group of this cover has order two.  The rank of its action
in the selected four-dimensional representation does not change that
fact.  Consequently every cyclic quotient supplier has only the paired
nonzero branch exponents \(m/2\), or is unramified for odd \(m\).
The complete supplier proof D-G34-04--05 and the full-centralizer Brauer
transfer D-G34-06 therefore apply.  In particular all primitive embeddings
have balanced even rank and receive nonzero algebraic determinant classes;
the faithful-character vanishing argument repairs the original-center
and exceptional-divisor step; and the full character group includes the
orthogonal parity coordinates.  The natural representation has a rational
one-copy factor by its displayed Schur-index-one realization.  The same
algebraic contraction as D-G24-04 gives the entire Weil plane on \(A\).
No assumption that the local involutions are all reflections is needed
in this supplier or transfer argument.

Define
\[
\begin{gathered}
 G=2304,\quad g=4609,\quad
 R=(2G+1)G!G^G,\quad L=2gR,\\
 S_0=G^{2g}(g!)^2,\quad S_B=(L/2)!S_0^R,\quad
 k=L/2+3,\\
 B_{0,2}=(40G)^k\binom{k}{3}S_B\,6!G^6,
 \qquad N_2=4G^6\Delta,\\
                 \boxed{D_{6,2}(\Delta)=N_2^6B_{0,2}.}
\end{gathered}
\]
Then every fibre of a connected fixed split \(\mathcal O_K\)-stable
Weil sixfold datum of polarization volume \(\Delta\) has an effective
integral three-cycle of degree at most \(D_{6,2}(\Delta)\) with nonzero
Weil projection.

\begin{proof}
The arithmetic supplier mass, divisor spanning, and moving arguments in
D-G24-05 are independent of the imaginary quadratic field, once the
source integral action, projector denominator, and actual dominance have
been supplied.  Here the projector denominator is \(G\), and
\eqref{eq:gf4-source-delta} replaces \(\Delta_A\mid(2G)^6\).
The coefficient mass is bounded by \(R\), the individual supplier mass
by \(S_0\), and the product mass by \(S_B\), exactly as in the checked
SQ3--SQ5 of \cite{G34:Hecke}.  The degree-six contracting tensor has
divisor degree \(c=\ell/2+3\le k\).  The effective mixed-intersection
estimate is
\[
 (40G)^c\binom c3\operatorname{mass}_{H_B}(Y_B)
                         \int_A\theta_A^6\le B_{0,2}.
\]
It gives a finite spanning collection of effective integral detector
three-cycles on the source; extracting a nonzero projected member incurs
no bound on rational or complex coefficients.

For the fixed target alternating form use
\[
                         h(x,y)=E(x,ty)+tE(x,y).
\]
This integral Hermitian form recovers \(E\) as the coefficient of \(t\).
Its determinant in signature \((3,3)\) is \(-\Delta\).  The ring
\(\mathbb Z[\sqrt{-2}]\) is norm Euclidean: approximation in its
rectangular lattice gives real and imaginary errors at most \(1/2\)
and \(\sqrt2/2\), with total norm at most \(3/4<1\).  It is a PID,
so the split-lattice comparison of D-G24-05 applies verbatim.  Both
directions have denominator dividing \(2\Delta\), and the source-to-target
comparison is cleared by \(4\Delta_A\Delta\), hence by \(N_2\).
The resulting isogeny pulls target polarization back to \(N_2^2\theta_A\)
and multiplies three-cycle degree by \(N_2^6\).  D-GF4-03 supplies the
nonempty analytic open image.  The accepted bounded-cycle closedness
then proves the displayed uniform bound on the whole fixed component.
\end{proof}

For the fixed order \(\mathcal O_f=\mathbb Z+f\mathcal O_K\), the
lattice \(f\mathcal O_K\Lambda\) has index at most \(f^{12}\), and its
pullback polarization has volume at most \(f^{12}\Delta\).  The same
projection-formula argument as D-G24-05 therefore gives
\[
                       \boxed{\delta_W\le D_{6,2}(f^{12}\Delta).}
\]
The project's accepted degree-to-division theorem makes this an explicit
bounded-division estimate for every fixed rational Weil class.  The result
includes the principal balanced \(E^6\) datum over \(\mathbb Q(\sqrt{-2})\).
It proves the split component and does not assert that nonsplit sixfold
determinants over this field have been covered.
